\section{Síntesis y ampliación}

La contribución más importante de esta clase no es afirmar que la IA ha reemplazado a los científicos, sino proporcionar un marco para \textbf{organizar modelos fundamentales en sistemas de colaboración científica}. Los modelos aportan conocimiento amplio y razonamiento; los Agentes asignan tareas de búsqueda; los torneos (tournaments) establecen prioridades relativas; las herramientas (tools) conectan el texto con cálculos especializados; la memoria preserva el estado entre rondas; y los humanos se encargan de los objetivos, las restricciones, los experimentos, la ética y la interpretación final.

\subsection{Comprimir toda la clase en un ciclo cerrado}

Los capítulos anteriores discutieron por separado la búsqueda, el ordenamiento, los experimentos y la seguridad; aquí se vuelven a conectar formando un ciclo con entrada, retroalimentación y condiciones de parada. Al leer esta cadena, no interprete las flechas como una tubería lineal única: si un experimento falla, regresa al Objetivo (Goal) y a la Generación (Generate); los expertos también pueden modificar restricciones en cualquier etapa, pausar la ejecución o solicitar evidencia adicional. El flujo de trabajo del científico colaborador con IA (AI co-scientist) se puede resumir como:
$$
\text{Goal}
\rightarrow \text{Generate}
\rightarrow \text{Debate}
\rightarrow \text{Rank}
\rightarrow \text{Evolve}
\rightarrow \text{Test}
\rightarrow \text{Update}.
$$
Donde:
\begin{itemize}
  \item El Objetivo (Goal) lo define el científico y puede reestructurarse durante el proceso;
  \item La Generación (Generate) y la Evolución (Evolve) amplían el espacio de candidatos;
  \item El Debate (Debate) y el Ranking (Rank) exponen debilidades y asignan atención interna;
  \item La Prueba (Test) conecta los candidatos con herramientas, expertos y experimentos;
  \item La Actualización (Update) modifica simultáneamente los candidatos, el grafo de evidencia, los registros de fallos y las preguntas para la siguiente ronda.
\end{itemize}

\begin{importantbox}{Las diez conclusiones centrales de esta clase}
\begin{enumerate}
  \item La velocidad del descubrimiento científico debe medirse por el conocimiento verificado, no por la cantidad de texto generado.
  \item La generación abierta de hipótesis requiere un proceso sistemático; no basta con una sola instrucción (prompt).
  \item AlphaFold demuestra que las tareas especializadas pueden comprimirse enormemente, pero eso no equivale a inteligencia científica general (general scientific intelligence).
  \item El valor de los sistemas multi-agente proviene de la división del trabajo heterogénea, la comparación y la retroalimentación, no de replicar respuestas idénticas.
  \item Elo y el modelo como juez (model-as-judge) solo proporcionan priorización interna; no pueden sustituir a la evidencia real.
  \item La reproducción oculta de resultados (AMR) no debe reescribirse como descubrimiento prospectivo.
  \item AML, organoides, predicción estructural y resultados no publicados en clase pertenecen a capas de evidencia distintas.
  \item Lo crucial en la validación prospectiva es congelar las predicciones y los métricas antes del experimento.
  \item Los informes extensos deben conservar el origen (provenance), contraejemplos, resultados negativos y metadatos de recursos.
  \item Tanto la participación del científico en el bucle (Scientist-in-the-loop) como la seguridad por capas (layered safety) deben贯穿 todo el proceso de ejecución, no añadirse al final con firmas póstumas.
\end{enumerate}
\end{importantbox}

\subsection{De throughput de hipótesis a throughput de validación}

Las limitaciones de los sistemas futuros probablemente no sean la incapacidad de generar hipótesis, sino la falta de experimentos suficientes, expertos y capacidad de revisión por pares. La optimización de ingeniería debe cambiar de "cuántos candidatos por hora" a "cuántos mecanismos se descartan y cuánta información se obtiene por unidad de costo de validación". Esto implica que los resultados negativos, el calibrado y las recomendaciones de parada tienen el mismo valor que los resultados positivos.

\begin{knowledgebox}{Lista de decisiones antes del despliegue}
¿Es falsable la pregunta? ¿Están congelados y rastreables los datos de entrada? ¿La división de tareas entre Agentes es realmente heterogénea? ¿Está pública la rúbrica del juez (judge rubric)? ¿Se registran todos los candidatos y resultados negativos? ¿Se congelan las predicciones antes del experimento? ¿Distingue el informe entre hechos, inferencias y recomendaciones? ¿Se publica el intervalo de cómputo (compute)? ¿Están minimizadas las permisos de las herramientas? ¿Hay un humano capaz de pausar, reestructurar y rechazar objetivos durante la ejecución?
\end{knowledgebox}

\subsection{Preguntas de investigación aún abiertas}

Primero, cómo evaluar la verdadera novedad sin recompensar estilos de escritura familiares para los modelos. Segundo, cómo permitir que múltiples Agentes compartan hechos manteniendo la diversidad de hipótesis. Tercero, cómo aprender del valor experimental sin simplificar la retroalimentación costosa del mundo real en sustitutos (proxy). Cuarto, cómo proporcionar un origen (provenance) a nivel de detalle y pruebas de falsificación ejecutables para los informes científicos. Quinto, cómo hacer que los sistemas de seguridad comprendan intenciones peligrosas que emergen solo después de combinaciones largas y cruzadas entre Agentes.

\begin{warningbox}{Conclusiones que no se pueden extraer de esta clase}
No se puede concluir que el científico colaborador con IA ya pueda realizar la ciencia por sí mismo; no se puede concluir que todos los casos mostrados estén publicados o sean utilizables clínicamente; no se puede concluir que más cómputo en tiempo de prueba (test-time compute) garantice inevitablemente un avance; no se puede concluir que un umbral de monitoreo del 10\% sea suficiente para asegurar la seguridad; ni tampoco se debe rellenar contenido no grabado de AMIE basándose en tarjetas de transición al final.
\end{warningbox}

\subsection{Lecturas adicionales}

Se recomienda primero leer la versión v1 de `Towards an AI co-scientist` de arXiv publicada antes de la clase, prestando especial atención a la arquitectura, el escalado en tiempo de prueba (test-time scaling), AML, AMR y el apéndice OCT4; luego leer la versión v1 de AMR de bioRxiv para verificar la línea temporal de la reproducción de resultados ocultos y las comparaciones entre candidatos; posteriormente consultar los dos artículos de bioRxiv v1 sobre fibrosis hepática y ensamblaje vegetal, comparando los niveles de validación de organoides y biología estructural. Al leer versiones posteriores en Nature, debe quedar claro que se publicaron después de la fecha de la clase y no deben interpretarse como evidencia del día de la clase.

\begin{knowledgebox}{Instantánea de la clase}
\raggedright
Esta clase tuvo lugar el 2026-05-14. Las apuntes utilizan cuatro versiones ya públicas en ese momento: AI co-scientist arXiv v1 (2502.18864), AMR bioRxiv v1 (2025.02.19.639094), liver-fibrosis bioRxiv v1 (2025.04.29.651320) y plant-assembly bioRxiv v1 (2026.05.03.722499). Las versiones de arXiv v2 del 2026-06-29, las versiones posteriores en Nature y otros materiales posteriores a la clase no constituyen hechos de la clase.
\end{knowledgebox}

\enlargethispage{3\baselineskip}
Finalmente, el valor del científico colaborador con IA no radica en producir una respuesta que parezca un "artículo", sino en convertir preguntas abiertas en un conjunto de candidatos comparables, refutables y enrutable hacia los expertos y experimentos adecuados. Mientras la validación, la responsabilidad y la seguridad permanezcan dentro del ciclo cerrado, los Agentes acelerarán la ciencia, no la certeza no verificada.

\end{document}
